%%%%%%%%%%%%%%%%%%%%%%%%%%%%%%%%%
%
%       EXERCÍCIO
%
%%%%%%%%%%%%%%%%%%%%%%%%%%%%%%%%%

\definevalorquestao

\begin{exercicioBanco}[\valorquestao]
Uma empresa de telefonia estima que o lucro diário \(L\) (em reais) obtido com determinado plano é dado pela expressão
\[
L(x)=30x-2x^2,
\]
em que \(x\) representa o número de dezenas de clientes que aderiram ao plano naquele dia.

Determine o valor do lucro \(L\), em reais, quando \(x=5\).

% Define as alternativas
\newcommand{\alternativas}{%
    \begin{center}
        \begin{tabularx}{\textwidth}{XXXXX}
            (a) \(100\). &
            (b) \(90\). &
            (c) \(80\). &
            (d) \(75\). &
            (e) \(50\).
        \end{tabularx}
    \end{center}
}

% Define a resposta correta
\newcommand{\resposta}{A}

% Lógica condicional para exibição
\ifdefstring{\atividade}{lista}{%
    \alternativas
    \vspace{0.5em}
    
    \noindent\textbf{Resposta:} letra \textbf{\resposta}.
}{%
    \ifdefstring{\modo}{objetivo}{%
        \alternativas
    }{%
        % modo = discursiva → não mostra alternativas
    }
}
\end{exercicioBanco}

